% =====================================================================
\section{Glossary}
{\small
\begin{description}[leftmargin=1.2em,itemsep=1pt,parsep=0pt,font=\normalfont\bfseries]
\item[AMPA receptor] the fast glutamate receptor
\item[ATP] the molecule in which a cell stores energy
\item[Astrocyte] a supporting cell of the brain; supplies neurons, including with the raw material for glutathione
\item[Axon] the output process of a neuron
\item[BDNF] brain-derived neurotrophic factor; a support signal released by active cells
\item[Catatonia] a state with severe motor disturbance; a separate diagnosis in ICD-11
\item[Cell culture] cells grown outside an organism
\item[Dendrite] an input process of a neuron
\item[Disinhibition] a rise in activity due to weakened inhibition
\item[Dopamine] a transmitter associated with the appraisal of the significance of events
\item[Enzyme] a protein that accelerates a chemical reaction
\item[Expression] production of a protein from a gene
\item[GABA] the main inhibitory transmitter
\item[GAD67] the enzyme that produces GABA
\item[Gamma rhythm] oscillations of cortical activity at 30--80\,Hz
\item[Glutamate] the main excitatory transmitter
\item[Glutathione] the principal consumable antioxidant of the cell
\item[Hippocampus] a brain region linked with memory; in Grace's scheme, the source of the excess signal
\item[ICD-11] the latest revision of the International Classification of Diseases
\item[Interneuron] an inhibitory neuron of the cortex
\item[Ketamine] a substance that blocks the NMDA receptor
\item[Master clock] the brain region above the crossing of the optic nerves (the suprachiasmatic nucleus), which synchronises the body's daily rhythms
\item[Membrane] the envelope of a cell
\item[Mitochondrion] the cell's power station; produces oxidants as a by-product
\item[mRNA] the working copy of a gene from which a protein is built
\item[NMDA receptor] the slow glutamate receptor; passes calcium and serves as an activity sensor
\item[Nox2] an enzyme that produces oxidants deliberately
\item[Oxidant] a molecule that takes electrons from other molecules
\item[Oxidative stress] an excess of oxidants over the capacity to clear them
\item[Parvalbumin (PV)] a calcium-binding protein; the marker of the fast inhibitory cells
\item[Perineuronal net] a shell of extracellular material around a PV cell
\item[Peroxiredoxin] the principal peroxide-scavenging enzyme
\item[Prodrome] the period before the first episode of psychosis
\item[Protein] a chain of amino acids assembled from a gene; the cell's main working material
\item[Psychosis] a state with hallucinations and delusions
\item[Pyramidal cell] an excitatory neuron of the cortex
\item[Receptor] a protein that receives a signal; in neurons, a gate for ions opened by a transmitter
\item[Serotonin] a transmitter governing the overall operating regime of the cortex; hallucinogens act on its receptors
\item[Sulfiredoxin] the enzyme that repairs hyperoxidised peroxiredoxin
\item[Synapse] the point of contact between two neurons
\item[Transmitter] a molecule that passes a signal across a synapse
\item[TrkB] the receptor that receives the BDNF signal
\end{description}}
